\documentclass[12pt]{article}

\usepackage[a4paper, margin=2.5cm]{geometry}
\usepackage{amsmath}
\usepackage{amssymb}
\usepackage{amsthm}
\usepackage[colorlinks=true, linkcolor=blue, citecolor=blue, urlcolor=blue]{hyperref}
\usepackage{booktabs}
\usepackage{natbib}

\newtheorem{theorem}{Theorem}
\newtheorem{proposition}[theorem]{Proposition}
\newtheorem{definition}[theorem]{Definition}
\newtheorem{remark}[theorem]{Remark}
\newtheorem{corollary}[theorem]{Corollary}
\newtheorem{conjecture}[theorem]{Conjecture}

\newcommand{\Mcal}{\mathcal{M}}
\newcommand{\Acal}{\mathcal{A}}
\newcommand{\phig}{\varphi}
\newcommand{\ds}{d_s}
\newcommand{\dseff}{d_s^{\mathrm{eff}}}

\title{Lepton Generations and the Missing Mass Operator\\in Deterministic Field Closure}
\author{%
  Lee Smart\\[4pt]
  \textit{Institute of Vibrational Field Dynamics}\\[2pt]
  \texttt{contact@vibrationalfielddynamics.org}\\[2pt]
  \texttt{@vfd\_org}%
}
\date{April 2026}

\begin{document}

\maketitle

\noindent\textbf{Status:} Preprint (not peer-reviewed)

\begin{abstract}
The three-order mass law of Paper~I~\citep{SmartI} recovers the proton-to-electron mass ratio at the $2 \times 10^{-4}$ level but cannot accommodate heavier lepton masses under shell-support extension. We prove a formal no-go for contiguous boundary-starting shell supports: the combinatorial invariant $C = k! - k(k{+}1)/2 - 1$ grows factorially, and no support of the form $S = \{1, 2, \ldots, k\}$ produces mass ratios in the muon or tau range. This establishes that lepton generations require a mass contribution independent of shell support.

The missing ingredient is a spectral layer supplying the exponent used in the winding ansatz. Within the present winding ansatz, the operator $f(w) = \phig^5 (w{-}1)^{1/\phig}$ requires an effective spectral dimension identified with $\ds \approx \phig$, which the base formalism does not supply. The F4 Fibonacci hierarchical construction~\citep{SmartIII} provides numerical evidence that $\dseff$ converges toward $\phig$ from above ($\dseff = 1.635$ at $V = 4960$, extrapolating to $1.628$ with $R^2 = 0.996$). Under the conjectural identification $\ds = \phig$, the winding operator recovers the muon-to-electron ratio to $0.5\%$ and the tau-to-electron ratio to $3.7\%$ with zero fitted continuous parameters. The spectral identification $\ds = \phig$ is conjectured, not derived; the coefficient $\phig^5$ is structurally selected, not derived; both remain open for future work.
\end{abstract}

%==================================================================
\section{Introduction}\label{sec:intro}
%==================================================================

Paper~I~\citep{SmartI} establishes a three-order mass law for the baryon-lepton hierarchy within the deterministic closure framework, recovering the proton-to-electron mass ratio at the $\sim\!2 \times 10^{-4}$ level with zero fitted continuous parameters. That law, however, cannot reproduce heavier lepton masses. The present paper shows that this limitation is structural, provable as a no-go theorem (Section~\ref{sec:nogo}), and identifies the nature of the missing mechanism.

Within the present framework, the required additional mechanism is a spectral layer supporting a winding-dependent mass contribution. Under the shell-support graph used in the base formalism, the effective spectral dimension is path-like ($\ds = 1$), not the value $\ds = \phig$ required by the winding ansatz.

The present paper has a narrower scope than a full theory of lepton masses. Its main derived result is a no-go theorem for contiguous shell-support extension. Its main positive result is conditional: if the relevant effective spectral dimension is identified with $\phig$, then the winding ansatz yields the muon and tau mass ratios with the accuracies reported below. The F4 analysis~\citep{SmartIII} is presented as numerical evidence supporting, but not proving, that identification. An additional spectral connection --- the eigenvalue $\lambda = 15$ of the 600-cell Laplacian, equal to $\Delta C$ --- links the mass law's leading term directly to the polytope spectrum (Section~\ref{sec:600cell}).

Section~\ref{sec:600cell} presents the 600-cell spectral connection. Section~\ref{sec:F4} shows that the F4 Fibonacci hierarchical construction converges toward $\dseff \approx \phig$, providing the primary numerical evidence for the spectral identification. Sections~\ref{sec:operator}--\ref{sec:predictions} present the winding operator and its mass predictions. Section~\ref{sec:dependency} gives the full dependency chain, and Section~\ref{sec:status} the claim status table.

%==================================================================
\section{No-Go Theorem for Shell-Extension Leptons}\label{sec:nogo}
%==================================================================

\textsc{Status: derived.}

\begin{proposition}[No-Go for Contiguous Boundary-Starting Shell Extensions]\label{prop:nogo}
Under the three-order mass law of~\citep{SmartI} with assignment rules R1--R5, no boundary-starting shell support $S = \{1, 2, \ldots, k\}$ produces a mass ratio $m(S)/m_e$ in the interval $[100, 4000]$.
\end{proposition}

\begin{proof}
For $S = \{1, \ldots, k\}$, the combinatorial invariant is $C(S) = k! - k(k{+}1)/2 - 1$. The three-order exponent relative to the electron is:
\[
\Delta\mathcal{E}(k) = [C(\{1,\ldots,k\}) - C(\{1\})] + \Delta g_1 + \Delta g_2
\]

Direct computation gives:

\begin{table}[ht]
\centering
\caption{Mass ratios under shell-support extension from boundary.}
\label{tab:nogo}
\begin{tabular}{@{}ccccc@{}}
\toprule
$k$ & $C$ & $\Delta\mathcal{E}$ & Computed ratio & Observed target \\
\midrule
1 & $-1$ & 0 & 1.00 & (electron) \\
2 & $-2$ & $-0.50$ & 0.79 & muon: 206.8 \\
3 & $-1$ & $+0.62$ & 1.35 & tau: 3477 \\
4 & $13$ & $+14.70$ & 1182 & --- \\
5 & $104$ & $+105.76$ & $>10^{21}$ & --- \\
\bottomrule
\end{tabular}
\end{table}

The function $k! - k(k{+}1)/2 - 1$ grows super-exponentially: $C$ jumps from $-1$ (at $k = 3$) to $13$ (at $k = 4$), producing a ratio gap from $\sim\!1$ to $\sim\!1200$ with no intermediate values. The muon ($\sim\!207$) and tau ($\sim\!3477$) mass ratios lie in this gap. No integer $k$ reaches the interval $[100, 4000]$. \qed
\end{proof}

\begin{corollary}
Lepton generations require a mass-lifting mechanism independent of shell-support size.
\end{corollary}

%==================================================================
\section{The Spectral Gap}\label{sec:gap}
%==================================================================

\textsc{Status: derived obstruction.}

The no-go theorem constrains the missing mechanism. It must:
\begin{enumerate}
    \item generate mass ratios of $\sim\!10^2$--$10^3$ for boundary leptons,
    \item not modify the proton-to-electron ratio (which uses winding $w = 1$ for both particles),
    \item be discrete (to preserve the discrete particle spectrum),
    \item have a structural origin within the closure framework.
\end{enumerate}

Within the present rule set, the only remaining discrete degree of freedom available for a boundary lepton (shell $\{1\}$, fixed by R1--R2) is the \textit{winding number} $w$. A winding-dependent boundary excitation operator $f(w)$ with $f(1) = 0$ satisfies constraint~(2) by construction.

Within the present winding ansatz, the $w$-th excitation is taken to scale as $(w - 1)^{d/(d+1)}$ in an effective dimension $d$. Recovering the muon mass motivates the condition:
\[
\frac{d}{d+1} = \frac{1}{\phig}
\]

This equation has the unique solution $d = \phig$, via the defining property of the golden ratio: $\phig/({\phig+1}) = 1/\phig$. The identity is exact.

The \textit{obstruction} is that under the shell-support graph used in the base formalism, the effective spectral dimension is path-like and evaluates to $\ds = 1$, not to $\ds = \phig$. The exponent $1/\phig$ is therefore conditionally available --- it follows rigorously from $\ds = \phig$ --- but the spectral dimension itself must come from outside the base shell structure.

%==================================================================
\section{Canonical Spectral Structure of the 600-Cell}\label{sec:600cell}
%==================================================================

\textsc{Status: theorem (eigenvalue connection); retracted (spectral dimension claim).}

\subsection{The 600-cell vertex graph}

The 600-cell is a regular four-dimensional polytope with 120 vertices, 720 edges, 1200 triangular faces, and 600 tetrahedral cells. Its vertex adjacency graph $G_{600}$ has $|V| = 120$ vertices, each of degree 12. It is a canonical geometric object: one of the six regular 4-polytopes, uniquely determined by its Schl\"afli symbol $\{3,3,5\}$.

\subsection{Eigenvalue connection: $\lambda = 15 = \Delta C$}

The graph Laplacian of $G_{600}$ has a discrete spectrum. Direct computation yields the eigenvalue $\lambda = 15$ with multiplicity 16. The closure invariant difference between proton and electron is $\Delta C = 15$. The mass law's leading-order exponent is therefore an eigenvalue of the Laplacian of the polytope that defines the framework's geometry.

This is a theorem-level connection --- verifiable by diagonalising the $120 \times 120$ Laplacian matrix. No windowing, fitting, or numerical estimation is involved.

\subsection{Retraction of the spectral dimension claim}

An earlier version of this paper reported the effective spectral dimension of the 600-cell vertex graph as $\dseff \approx 1.637$--$1.642$, near $\phig$. A more careful analysis with finer time resolution shows that $\dseff(t)$ does not plateau near $\phig$; the earlier value was an artifact of coarse windowing across mixed spectral regimes. The $\dseff \approx \phig$ claim for the 600-cell vertex graph is retracted.

The F4 Fibonacci hierarchical construction~\citep{SmartIII}, which converges toward $\dseff \approx 1.63$ at $V = 4960$ with proper plateau detection across multiple scales, remains independent of this retraction and provides the primary numerical evidence for the spectral identification $\ds \approx \phig$.

%==================================================================
\section{Interpretation}\label{sec:interpretation}
%==================================================================

The 600-cell connects to the mass law through the Laplacian eigenvalue $\lambda = 15 = \Delta C$, not through its effective spectral dimension. This connection is a different --- and arguably more direct --- link than the spectral-dimension route: the mass exponent's leading term is a specific point in the polytope spectrum, verifiable by exact computation.

%==================================================================
\section{F4 Fibonacci Hierarchical Construction}\label{sec:F4}
%==================================================================

\textsc{Status: primary numerical evidence for $\ds \approx \phig$.}

A companion paper~\citep{SmartIII} introduces the F4 Fibonacci hierarchical extension of the $\phig$-manifold: a construction in which shell~$k$ contains $F(k{+}2)$ sub-cells, where $F$ is the Fibonacci sequence. The F4 graph is a constructed hierarchical model rather than a canonical polytope graph; its role is to provide scalable numerical evidence that a $\phig$-structured geometry can produce an effective spectral dimension converging toward $\phig$.

Key results from~\citep{SmartIII}:
\begin{itemize}
    \item The effective spectral dimension of the F4 family converges from above toward $\phig$ as the base size $b$ increases: $\dseff^{HK} = 1.635$ at $V = 4960$.
    \item Empirical $1/V$ extrapolation gives $\ds(\infty) = 1.628$ ($R^2 = 0.996$).
    \item Among all constructions tested (binary tree, ternary tree, non-Fibonacci hierarchical, 1D chain), F4 is the only one among the tested control families whose effective spectral dimension trends toward $\phig$.
    \item The random-walk estimate is robust to $20\%$ edge-weight perturbation ($\Delta \dseff \approx 0.001$).
\end{itemize}

The F4 construction is the primary numerical evidence for $\dseff \to \phig$: (i) it shows monotone convergence from above toward $\phig$; (ii) among all tested control families, it is the only one whose effective spectral dimension trends toward $\phig$; and (iii) it demonstrates that an effective spectral dimension near $\phig$ requires specific hierarchical structure.

%==================================================================
\section{The Winding Excitation Operator}\label{sec:operator}
%==================================================================

\subsection{Operator form}

The winding contribution to the mass exponent is:
\begin{equation}\label{eq:fw}
    f(w) = \phig^N \cdot (w - 1)^{1/\phig}
\end{equation}
where $N$ is the manifold shell count and $\phig = (1+\sqrt{5})/2$. This satisfies $f(1) = 0$ by construction.

\subsection{Status of the exponent $1/\phig$}

\textsc{Status: conditional} on $\ds \approx \phig$.

The exponent $1/\phig$ follows within the present spectral-scaling ansatz $d/(d+1)$ under the identification $d = \ds = \phig$. The identity $\phig/(\phig+1) = 1/\phig$ is exact (it is the defining property of the golden ratio). The exponent is therefore conditionally derived: it follows rigorously from the spectral-dimension identification, which is itself conjectured.

\subsection{Status of the coefficient $\phig^N$}

\textsc{Status: structurally selected.}

Within the restricted family $\phig^k$ with integer $k$, the choice $k = 5$ is the only member within $1\%$ of the coefficient required by the muon mass:

\begin{table}[ht]
\centering
\caption{Coefficient candidates in the $\phig^k$ family.}
\label{tab:coefficient}
\begin{tabular}{@{}cccc@{}}
\toprule
$k$ & $\phig^k$ & Muon error & Status \\
\midrule
4 & 6.854 & 38.1\% & Too small \\
\textbf{5} & \textbf{11.090} & \textbf{0.097\%} & \textbf{Only near-match in tested $\phig^k$ family} \\
6 & 17.944 & 62.0\% & Too large \\
\bottomrule
\end{tabular}
\end{table}

The interpretation of $N = 5$ as the total manifold shell count gives a structural meaning: the boundary excitation scale equals the total depth of the $\phig$-manifold. This is an interpretation within the framework, not an independent derivation.

\textsc{Status of $N = 5$: postulated.} The shell count is an input of the framework, not derived from it.

%==================================================================
\section{Mass Predictions}\label{sec:predictions}
%==================================================================

Lepton generations are reinterpreted as winding excitations of the boundary closure class:
\begin{align}
    \text{electron:} \quad & \{1\},\; w = 1 \quad \text{(ground state)} \\
    \text{muon:} \quad & \{1\},\; w = 2 \quad \text{(first excitation)} \\
    \text{tau:} \quad & \{1\},\; w = 3 \quad \text{(second excitation)}
\end{align}

With $f(w) = \phig^5 (w{-}1)^{1/\phig}$ and zero fitted continuous parameters:

\begin{table}[ht]
\centering
\caption{Lepton mass ratios under the winding operator. These results are conditional on the spectral identification $\ds \approx \phig$.}
\label{tab:predictions}
\begin{tabular}{@{}llccl@{}}
\toprule
Particle & Winding & Computed ratio & Observed & Agreement \\
\midrule
Electron & $w = 1$ & 1 (reference) & --- & --- \\
Muon & $w = 2$ & $\phig^{\phig^5} \approx 207.8$ & 206.8 & $0.5\%$ \\
Tau & $w = 3$ & $\phig^{\phig^5 \cdot 2^{1/\phig}} \approx 3607$ & 3477 & $3.7\%$ \\
\bottomrule
\end{tabular}
\end{table}

The muon value uses $f(2) = \phig^5$. The tau value is a genuine test: $f(3) = \phig^5 \cdot 2^{1/\phig}$ is determined entirely by the operator form, not calibrated to the tau.

The proton-to-electron ratio is unaffected: both particles have $w = 1$, so $f(1) = 0$. The Paper~I result $m_p/m_e = \phig^{1265/81} \approx 1835.8$ is exactly preserved.

%==================================================================
\section{Extended Factorised Law}\label{sec:unified}
%==================================================================

The extended mass law, conditional on the winding ansatz and incorporating the Paper~I baryon-lepton law, takes the factorised form:
\begin{equation}\label{eq:unified}
    m(\Acal) = \Lambda \;\cdot\; \phig^{C(\Acal)} \;\cdot\; \phig^{|E|/|V|} \;\cdot\; \phig^{-\mathrm{Var}(\deg)\,|E|/|V|^2} \;\cdot\; \phig^{f(w)}
\end{equation}
where $|E|$ and $|V|$ are the edge and vertex counts of the closure graph, and $f(w)$ is defined by Eq.~\eqref{eq:fw}. In the present framework, each factor is assigned to a distinct structural sector and the leading-order law is taken to factorise across those sectors. For particles with $w = 1$ (electron, proton, neutron), the winding factor equals $\phig^0 = 1$ and the law reduces to Paper~I.

%==================================================================
\section{Dependency Chain}\label{sec:dependency}
%==================================================================

The logical structure of the lepton generation programme is:

\begin{center}
\begin{tabular}{@{}lcl@{}}
F4 Fibonacci construction & $\xrightarrow{\text{numerical convergence}}$ & $\dseff \to \phig$ \\[4pt]
$\dseff \approx \phig$ & $\xrightarrow{\text{conjectural identification}}$ & $\ds = \phig$ \\[4pt]
$\ds = \phig$ & $\xrightarrow{\text{exact identity}}$ & exponent $= 1/\phig$ \\[4pt]
exponent $+ \phig^5$ coefficient & $\xrightarrow{\text{winding operator}}$ & $f(w) = \phig^5 (w{-}1)^{1/\phig}$ \\[4pt]
$f(w)$ & $\xrightarrow{\text{evaluation}}$ & lepton mass ratios \\[8pt]
\multicolumn{3}{@{}l@{}}{\textit{Independent support:} 600-cell Laplacian eigenvalue $\lambda = 15 = \Delta C$ (theorem).} \\
\end{tabular}
\end{center}

Each arrow carries a stated epistemic status. The chain is as strong as its weakest link: the conjectural identification $\ds = \phig$.

%==================================================================
\section{Claim Status}\label{sec:status}
%==================================================================

\begin{table}[ht]
\centering
\caption{Claim status table.}
\label{tab:status}
\begin{tabular}{@{}lll@{}}
\toprule
Claim & Status & Basis \\
\midrule
Shell-extension no-go & \textsc{derived} & Exhaustive computation \\
Spectral gap obstruction & \textsc{derived} & Base shell-support graph gives path-like $\ds = 1$ \\
$\lambda = 15$ (600-cell eigenvalue) & \textsc{theorem} & Laplacian diagonalisation; $\lambda = \Delta C$ \\
$\ds = \phig$ & \textsc{conjectured} & Proximity; not proven \\
Exponent $1/\phig$ & \textsc{conditional} & Requires $\ds = \phig$ \\
Coefficient $\phig^5$ & \textsc{structurally selected} & Unique in $\phig^k$ family \\
Shell count $N = 5$ & \textsc{postulated} & Framework input \\
Full winding operator & \textsc{conditional} & Requires $\ds = \phig$ and $N = 5$ \\
Muon ratio ($0.5\%$) & \textsc{conditional} & Follows from operator \\
Tau ratio ($3.7\%$) & \textsc{conditional} & Genuine test (not calibrated) \\
F4 convergence toward $\phig$ & \textsc{numerical observation} & Companion paper~\citep{SmartIII} \\
\bottomrule
\end{tabular}
\end{table}

%==================================================================
\subsection*{What this paper does not claim}
%==================================================================

This paper does not claim an analytic proof that $\ds = \phig$, nor that the winding operator is fully derived from first principles. It does not derive the shell count $N = 5$ or the coefficient $\phig^5$. An earlier claim that the 600-cell vertex graph has effective spectral dimension $\dseff \approx 1.637$--$1.642$ near $\phig$ has been retracted (Section~\ref{sec:600cell}); the F4 construction provides the primary numerical evidence for $\ds \approx \phig$. The paper's main claims are: a derived shell-extension no-go theorem, the eigenvalue connection $\lambda = 15 = \Delta C$ on the 600-cell Laplacian, and a conditional winding operator supported by the F4 spectral evidence.

%==================================================================
\section{Open Problems}\label{sec:open}
%==================================================================

\begin{enumerate}
    \item \textbf{Exact proof of $\ds = \phig$.} The 600-cell observation ($\dseff \approx 1.64$) and the F4 extrapolation ($\ds(\infty) = 1.628$) are consistent with $\ds = \phig$ but do not prove it. An analytical derivation --- from the spectral zeta function, renormalisation group, or heat-kernel asymptotics --- remains the primary formal target.

    \item \textbf{Coefficient derivation.} The coefficient $\phig^5$ is structurally selected (unique in its family) but not derived from the closure geometry. Whether the F4 spectral recursion or the 600-cell structure provides an independent route to $\phig^5$ is unknown.

    \item \textbf{Canonical selection.} Why the 600-cell vertex geometry, rather than another polytope? The observation that only this specific vertex graph exhibits the required spectral behaviour is suggestive but not constitutive. A structural argument selecting the 600-cell from within the closure framework would strengthen the programme.

    \item \textbf{Tau mass correction.} The tau ratio has $3.7\%$ structural error. Whether a higher-order winding correction, analogous to the graph corrections in Paper~I, can reduce this is unknown.

    \item \textbf{Neutron-proton splitting.} A symmetry-sector correction $\phig^{\delta_\sigma}$ with $\delta_\sigma \approx 0.003$ is compatible with the architecture but not derived.

    \item \textbf{Extension to quarks and heavier hadrons.}
\end{enumerate}

%==================================================================
\section{Discussion}\label{sec:discussion}
%==================================================================

The central structural result is the no-go theorem: shell-support extension cannot produce lepton generations. This identifies the class of the missing mechanism (spectral, winding-dependent) without determining it uniquely.

The 600-cell connects to the mass law through the Laplacian eigenvalue $\lambda = 15 = \Delta C$ --- a theorem-level spectral embedding of the closure invariant. The F4 Fibonacci hierarchical construction provides the primary numerical evidence that a $\phig$-structured geometry can produce $\dseff \to \phig$, supporting the conditional winding operator. An earlier claim that the 600-cell vertex graph itself has $\dseff \approx \phig$ has been retracted; the F4 evidence stands independently.

The winding operator $f(w) = \phig^5 (w{-}1)^{1/\phig}$, conditional on the spectral identification, recovers three charged lepton masses from a single functional form with zero fitted continuous parameters. The agreement is structural rather than metrological: the muon and tau masses are known to much higher precision than $0.5\%$ and $3.7\%$ respectively.

Whether the spectral identification $\ds = \phig$ can be established analytically is the primary remaining formal question. If it can, the exponent $1/\phig$ becomes derived rather than conditional, upgrading the status of the entire lepton programme.

%==================================================================
\section{Conclusion}\label{sec:conclusion}
%==================================================================

A formal no-go theorem establishes that lepton generations cannot arise from shell-support extension within the closure-class mass law. The no-go theorem identifies the need for a winding-dependent boundary excitation operator, and the form $f(w) = \phig^5 (w{-}1)^{1/\phig}$ is proposed as the minimal structurally compatible candidate within the present framework. The exponent $1/\phig$ follows within the spectral-scaling ansatz under the identification $\ds = \phig$; this identification is supported by the F4 Fibonacci hierarchical construction ($\dseff = 1.635$ at $V = 4960$, extrapolating to $1.628$), but has not been analytically proven. An earlier claim that the 600-cell vertex graph has $\dseff \approx \phig$ has been retracted; the 600-cell's connection to the mass law is instead through the Laplacian eigenvalue $\lambda = 15 = \Delta C$.

The operator recovers the muon-to-electron ratio to $0.5\%$ and the tau-to-electron ratio to $3.7\%$ with zero fitted continuous parameters. Together with the proton-to-electron law of Paper~I~\citep{SmartI}, the framework provides a constrained, conditional structural extension from the baryon-lepton hierarchy to the three charged leptons.

The programme's epistemic structure is explicit: the no-go and spectral gap are derived; the 600-cell spectral dimension is a numerical observation; the identification $\ds = \phig$ is conjectured; the winding operator and its mass predictions are conditional on that conjecture. Closing the spectral gap analytically is the primary remaining formal target.

%==================================================================
\bibliographystyle{plainnat}
\bibliography{references}

\end{document}
